% RS_Arithmetic_From_Law_Of_Logic.tex
%
% "Arithmetic from the Law of Logic:
%  The Continuous Positive-Ratio Realization."
%
% Companion to:
%   Washburn, "A Functional-Equation Encoding of Logical Consistency
%   on Continuous Positive-Ratio Comparisons," April 2026
%   (Logic_Functional_Equation.tex / Mathematical_Foundation/).
%
% Cross-realization invariance is the subject of the meta-theorem
% companion paper RS_Universal_Forcing.tex; this paper supplies the
% concrete continuous-ratio realization on which the meta-theorem
% rests. The original "Part I" framing has been retired because the
% meta-theorem subsumes subsequent installments and worked examples
% of other realizations are already in RS_Universal_Forcing.tex.
%
% Formalization details are collected in the implementation appendix.

\documentclass[11pt,a4paper]{article}

\usepackage[T1]{fontenc}
\usepackage{lmodern}
\usepackage{amsmath, amssymb, amsthm}
\usepackage[margin=1in]{geometry}
\usepackage{hyperref}
\usepackage{microtype}
\usepackage{enumitem}
\usepackage{listings}

\hypersetup{colorlinks=true, linkcolor=black, citecolor=blue, urlcolor=blue}

\newtheorem{theorem}{Theorem}
\newtheorem{lemma}[theorem]{Lemma}
\newtheorem{proposition}[theorem]{Proposition}
\newtheorem{corollary}[theorem]{Corollary}

\theoremstyle{definition}
\newtheorem{definition}[theorem]{Definition}

\theoremstyle{remark}
\newtheorem{remark}[theorem]{Remark}

\title{Arithmetic from the Law of Logic:\\
The Continuous Positive-Ratio Realization}

\author{Jonathan Washburn\thanks{Recognition Physics Institute, Austin, TX 78701, USA.\\
\texttt{jon@recognitionphysics.org}.}}

\date{April 2026}

\begin{document}

\maketitle

\begin{abstract}
The companion paper~\cite{Washburn2026LogicFE} proves a
\emph{conditional rigidity theorem} on the continuous positive-ratio
comparison setting: under the four classical Aristotelian conditions
on a comparison operator, plus continuity and finite pairwise
polynomial closure of the combiner, the derived cost function
satisfies the Recognition Composition Law and (under unit
log-curvature calibration) equals the canonical reciprocal cost
$J(x) = \tfrac{1}{2}(x + x^{-1}) - 1$. \emph{Conditional on that
rigidity theorem} (and on the inductive-type, quotient, and
completion machinery of the ambient meta-language), we exhibit the
standard Peano structure as the canonical orbit of any non-trivial
generator inside the multiplicative group $\mathbb{R}_{>0}$, lift the
construction through the integers, rationals, reals, and the complex
carrier, and identify each layer up to a transport equivalence with
the usual number tower. The Peano core carries no domain-specific
axiomatic content; the upper layers reuse standard algebraic and
completion machinery through explicit transport equivalences. The
companion meta-theorem~\cite{Washburn2026UniversalForcing}, the
\emph{Universal Forcing Theorem}, asserts that in every admissible
realization of the Law of Logic the same arithmetic structure is
canonically forced up to unique isomorphism; the present paper
supplies the concrete continuous-ratio realization on which that
meta-theorem rests. The downstream consequence is that the
elementary number tower used by Recognition Science is forced by a
single law, realised concretely here on continuous positive ratios
and invariant across all admissible realizations of that law, rather
than imported as a separate axiomatic commitment, conditional on
the named hypotheses recorded in Section~\ref{subsec:conditionality}.
\end{abstract}

\noindent\textbf{Keywords:} laws of logic; functional equation;
Recognition Composition Law; Peano arithmetic; Grothendieck completion;
Cauchy completion; formal mathematics.

\medskip

\noindent\textbf{MSC 2020:} 03A05, 03B22, 03F35, 39B22, 68V20.

\section{Introduction}

\paragraph{Setup.} The companion paper \cite{Washburn2026LogicFE}
identifies, on the continuous-ratio comparison setting, the unique
cost function compatible with an operational encoding of the four
classical Aristotelian conditions on a continuous comparison operator
$C : \mathbb{R}_{>0} \times \mathbb{R}_{>0} \to \mathbb{R}$. Given the
encoding (Identity, Non-Contradiction, Excluded Middle, Composition
Consistency) plus continuity and finite pairwise polynomial closure,
the derived cost $F(r) := C(r, 1)$ satisfies the Recognition
Composition Law family, and its calibrated representative on the
bilinear branch is the canonical reciprocal cost
$J(x) = \tfrac{1}{2}(x + x^{-1}) - 1$.

\paragraph{Question.} The companion paper explicitly does not claim a
derivation of arithmetic. We close that gap on the same restricted
domain. The Law of Logic does more than pin down the cost function on
this setting; it canonically picks out an inductive structure on the
multiplicative group of positive reals, and that inductive structure is
the natural-number structure.

\paragraph{The foundational theorem behind the program.} The thesis
is sharper than what any single setting can prove:

\begin{quote}
\textbf{Universal Forcing Theorem.} In every admissible realization
of the Law of Logic, the canonical forced inductive object is unique
up to unique isomorphism. Equivalently, the arithmetic content of
the Law of Logic is invariant under choice of realization carrier.
\end{quote}

The theorem cannot be discharged on a single setting. The present
paper supplies the continuous positive-ratio realization. The
companion paper~\cite{Washburn2026UniversalForcing} states and
proves the meta-theorem and surveys further realizations (discrete
propositional, modular cyclic, order-theoretic, Lawvere-style
categorical) as worked examples;
\S\ref{sec:program} of the present paper records the bridge.

\paragraph{Result of this paper.} Conditional on the
\cite{Washburn2026LogicFE} rigidity theorem, we recover Peano
arithmetic from the comparison-operator structure on continuous
positive ratios. We extend the recovery to the integers via
Grothendieck completion, to the rationals by field-of-fractions
construction, to the reals by Cauchy/Bourbaki completion, and to the
complex carrier by ordered pairs of recovered reals. At every
discrete stage the Peano axioms, the arithmetic identities, and the
order-theoretic content are theorems of the inductive structure
canonically picked out by the comparison operator, not posits. The
real and complex layers reuse standard completion and analytic
substrate machinery through explicit transport equivalences.

\subsection{Conditionality, stated up front}
\label{subsec:conditionality}

The result of this paper is conditional on three named layers,
listed here so the reader knows the dependency structure before
reading the body. Each layer is also discussed in its own context
later in the paper.

\begin{enumerate}[leftmargin=1.6em, label=(\arabic*)]
\item \textbf{The rigidity theorem of}~\cite{Washburn2026LogicFE}
for continuous positive-ratio comparisons. This is the only
setting-specific input. It is itself a \emph{conditional} rigidity
theorem under continuity and finite pairwise polynomial closure of
the combiner; the polynomial-degree-two regularity is essential, not
a removable hypothesis (the quartic-log counterexample
$C(x, y) = (\ln(x/y))^4$ in the companion paper establishes
sharpness). Our recovery of arithmetic is therefore doubly
conditional: it inherits the conditional structure of the rigidity
theorem and adds the universal-property characterisation of the
forced orbit.

\item \textbf{The inductive-type, quotient, and completion
machinery of the ambient meta-language.} The Peano core is
constructed as an inductive type with two constructors; the integers
are built by Grothendieck quotient; the rationals by
field-of-fractions quotient; the reals by Cauchy completion. None
of these constructions is novel; each is the standard meta-theoretic
machinery used in type-theoretic foundations and is a reasonable
baseline for any formal foundational result.

\item \textbf{Mathlib's analytic substrate} where the upper tower
is transported. The recovered $\mathbb{R}$ and $\mathbb{C}$ inherit
their analytic structure (limits, derivatives, integrals, holomorphy,
transcendental functions) by transport equivalence with Mathlib's
standard $\mathbb{R}$ and $\mathbb{C}$. This is honest reuse, not a
foundational claim; the recovered objects could in principle be
developed analytically from scratch, but doing so adds no
foundational content.
\end{enumerate}

The summary: the recovery of Peano $\mathbb{N}_{\mathrm{Logic}}$ is
\emph{doubly conditional} (on (1) the conditional rigidity theorem
of the companion paper, and on (2) the standard inductive machinery
of the meta-language). The recovery of $\mathbb{Z}_{\mathrm{Logic}}$
through $\mathbb{Q}_{\mathrm{Logic}}$ adds standard quotient
constructions on top of (1) and (2). The recovery of
$\mathbb{R}_{\mathrm{Logic}}$ and $\mathbb{C}_{\mathrm{Logic}}$ adds
the analytic transport in (3). All three layers are named
explicitly; the cover letter and any subsequent presentation of this
result should foreground the conditional structure honestly.

\paragraph{Distinguishing claim.} The natural numbers are not assumed
as an extra commitment of this paper. They are not invoked as a
primitive of set theory. They are picked out canonically as the orbit
of a single non-trivial element of $(\mathbb{R}_{>0}, \cdot)$, and
identified with the standard inductive object via a universal-property
characterization (\S\ref{sec:universal}). The standard textbook
presentation either takes $\mathbb{N}$ as a primitive (Peano
postulation) or builds it set-theoretically inside ZFC (von Neumann
ordinals, Zermelo construction). The present paper does neither.

\section{Background and Notation}
\label{sec:background}

We recall the encoding of \cite{Washburn2026LogicFE}. A
\emph{comparison operator} is a map $C : \mathbb{R}_{>0} \times
\mathbb{R}_{>0} \to \mathbb{R}$. The four classical Aristotelian
conditions are encoded as:

\begin{itemize}
\item \textbf{Identity}: $C(x, x) = 0$ for all $x > 0$.
\item \textbf{Non-Contradiction}: $C(x, y) = C(y, x)$ for all $x, y > 0$.
\item \textbf{Excluded Middle}: $C$ is jointly continuous and total on
the open positive quadrant.
\item \textbf{Composition Consistency}: there exists a polynomial
combiner $P$ of total degree at most two with
$F(xy) + F(x/y) = P(F(x), F(y))$, where $F(r) := C(r, 1)$.
\end{itemize}

We add two structural conditions:

\begin{itemize}
\item \textbf{Scale Invariance} (bridge): $C(\lambda x, \lambda y) = C(x, y)$
for all $\lambda, x, y > 0$.
\item \textbf{Non-Triviality}: there exists $x > 0$ with $F(x) \neq 0$.
\end{itemize}

A comparison operator satisfying these six conditions is said to
\emph{satisfy the Law of Logic}, abbreviated $C \in \mathrm{Logic}$.
The companion paper proves that, given $C \in \mathrm{Logic}$, the
derived cost $F$ satisfies the Recognition Composition Law and, under
calibration, equals $J$.

\section{The Inductive Structure Forced by a Non-Trivial Generator}
\label{sec:inductive}

\paragraph{Claim.} The Law of Logic forces an inductive structure
isomorphic to $\mathbb{N}$.

\paragraph{Construction.} Non-Triviality together with Identity gives
an element $\gamma \in \mathbb{R}_{>0}$ with $\gamma \neq 1$. (If
$F(x) \neq 0$ then $x \neq 1$, since $F(1) = C(1, 1) = 0$ by
Identity.) The orbit
\[
\mathcal{O}(\gamma) := \{1, \gamma, \gamma^2, \gamma^3, \dots\}
\]
is the smallest subset of $\mathbb{R}_{>0}$ containing 1 and closed
under multiplication by $\gamma$. The two-constructor presentation of
$\mathcal{O}(\gamma)$ matches the Peano signature: ``be at the
identity'' (constructor 1) or ``take one more step'' (constructor 2).

\begin{definition}
The recovered natural-number object is the inductive object with two
constructors: an identity constructor (nullary) and a step constructor
(unary, taking one recovered natural number to the next).
\end{definition}

The formal construction identifies this object with the orbit
structure via the embedding
\[
\mathrm{embed}_\gamma : \mathbb{N}_{\mathrm{Logic}} \to \mathbb{R}, \qquad
\mathrm{embed}_\gamma(0) = 1, \quad
\mathrm{embed}_\gamma(Sn) = \gamma \cdot \mathrm{embed}_\gamma(n).
\]

\section{Peano Axioms as Theorems}
\label{sec:peano}

The Peano axioms are theorems of the inductive type, not posits.

\begin{theorem}[Peano P1]
$0 \neq S(n)$ for every recovered natural number $n$.
\end{theorem}

\begin{theorem}[Peano P2]
The successor map $S$ is injective.
\end{theorem}

\begin{theorem}[Peano P3, induction]
For any predicate $P$ on recovered natural numbers, if $P(0)$ and
$\forall n,\, P(n) \to P(Sn)$, then $\forall n,\, P(n)$.
\end{theorem}

Each statement follows from inductive disjointness of constructors
and the recursor of the inductive object. The dependency audit is
recorded in Appendix~\ref{app:formal}.

\section{Addition, Multiplication, Order}
\label{sec:operations}

Addition is repeated successor; multiplication is repeated addition.
Both are defined by recursion on the second argument. The standard
identities follow by induction:

\begin{theorem}[Arithmetic identities]
The recovered natural-number object with the operations defined above
satisfies:
\[
0 + n = n, \quad n + 0 = n, \quad
n + m = m + n, \quad (a + b) + c = a + (b + c),
\]
\[
0 \cdot n = 0, \quad 1 \cdot n = n, \quad
a \cdot (b + c) = a \cdot b + a \cdot c.
\]
\end{theorem}

The \emph{order} on recovered natural numbers is the Peano definition:
$n \le m$ iff there exists $k$ with $n + k = m$. We prove
reflexivity, transitivity, antisymmetry, totality, decidability, and
well-foundedness.

\section{The Recovery Isomorphism}
\label{sec:recovery}

\begin{theorem}[Recovery]
There is an isomorphism $\mathbb{N}_{\mathrm{Logic}} \simeq \mathbb{N}$
under which addition, multiplication, and the strict order on
$\mathbb{N}_{\mathrm{Logic}}$ correspond to those on $\mathbb{N}$:
\[
\mathrm{toNat}(a + b) = \mathrm{toNat}(a) + \mathrm{toNat}(b),
\quad
\mathrm{toNat}(a \cdot b) = \mathrm{toNat}(a) \cdot \mathrm{toNat}(b),
\]
\[
a \le b \Leftrightarrow \mathrm{toNat}(a) \le \mathrm{toNat}(b),
\quad
a < b \Leftrightarrow \mathrm{toNat}(a) < \mathrm{toNat}(b).
\]
\end{theorem}

The isomorphism is established by mutual recursion and is compatible
with addition, multiplication, non-strict order, and strict order.

The axiom-dependency audit is given in Appendix~\ref{app:formal}.

\section{The Integers via Grothendieck Completion}
\label{sec:integers}

The integers are recovered as the Grothendieck completion of the
additive monoid $(\mathbb{N}_{\mathrm{Logic}}, +)$. Equivalence
\[
(a, b) \sim (c, d) \iff a + d = c + b
\]
on $\mathbb{N}_{\mathrm{Logic}} \times \mathbb{N}_{\mathrm{Logic}}$ has the standard
properties (reflexivity, symmetry, transitivity by additive
cancellation in $\mathbb{N}_{\mathrm{Logic}}$). The quotient is the
recovered integer object.

\begin{theorem}[Integer carrier recovery]
There is an equivalence $\mathbb{Z}_{\mathrm{Logic}} \simeq \mathbb{Z}$
under which the operations $+$, $-$, $\cdot$ defined by quotient lift
on recovered integers correspond, at the carrier level, to those on
$\mathbb{Z}$.
\end{theorem}

The construction gives the carrier-level equivalence and the ring
laws by transport through this equivalence.

\section{Rationals, Reals, and Complex Numbers}
\label{sec:tower}

The integer layer is not the endpoint. Once $\mathbb{Z}_{\mathrm{Logic}}$
has been recovered, the rationals are obtained by the standard field of
fractions construction:
\[
(a,b) \sim (c,d) \iff a d = c b,\qquad b,d\neq 0.
\]
The quotient is denoted $\mathbb{Q}_{\mathrm{Logic}}$, and the transport
map to $\mathbb{Q}$ preserves zero, one, negation, addition, and
multiplication.

\begin{theorem}[Rational recovery]
There is an equivalence
\[
\mathbb{Q}_{\mathrm{Logic}} \simeq \mathbb{Q}
\]
under which the field-of-fractions operations on
$\mathbb{Q}_{\mathrm{Logic}}$ correspond to the usual rational
operations.
\end{theorem}

The real layer is then obtained as the completion of the recovered
rationals. In the formalization this is implemented by transporting
$\mathbb{Q}_{\mathrm{Logic}}$ to Mathlib's rational line, using the
Bourbaki completion, and transporting back. Thus the theorem is not that
there is a new real line, but that the real line used by the analytic
parts of the theory sits at the end of the recovered number tower.

\begin{theorem}[Real recovery]
There is an equivalence
\[
\mathbb{R}_{\mathrm{Logic}} \simeq \mathbb{R}
\]
compatible with the transported algebraic operations and order.
\end{theorem}

The complex carrier is recovered as pairs of recovered reals:
\[
\mathbb{C}_{\mathrm{Logic}} :=
\mathbb{R}_{\mathrm{Logic}} \times \mathbb{R}_{\mathrm{Logic}}.
\]
Transport to Mathlib's $\mathbb{C}$ is coordinatewise, and addition,
negation, multiplication, inversion, and division are pulled back along
this equivalence.

\begin{theorem}[Complex carrier recovery]
There is a carrier equivalence
\[
\mathbb{C}_{\mathrm{Logic}} \simeq \mathbb{C}.
\]
\end{theorem}

We deliberately call this a carrier recovery. Complex analysis,
holomorphy, contour integration, and zeta-function arguments are still
performed in the established analytic substrate and are related back by
transport theorems. This avoids claiming that the present paper has
rederived all complex analysis from scratch.

\section{Universal Property of the Recovered Natural Numbers}
\label{sec:universal}

The reviewer-style objection ``you have just used the proof assistant's
inductive-type machinery to import $\mathbb{N}$ under a new name'' is
the right thing to confront here. The objection is partially true, in
the sense that any formal foundational paper must use the inductive
machinery of its metatheory; but the substantive content of this
section is that the comparison-operator structure picks out the
natural-number object up to unique isomorphism, with no choice in the
identification beyond the choice of generator.

\begin{theorem}[Generator-orbit initiality]
\label{thm:initiality}
Let $C$ satisfy the Law of Logic. For any non-trivial generator
$\gamma \in \mathbb{R}_{>0}$ extracted from
$\mathrm{NonTriviality}(C)$, the triple $(\mathcal{O}(\gamma), 1, \cdot
\gamma)$ is the initial Peano structure inside the comparison-operator
setting. Concretely, for every triple $(X, x_0, f)$ in the same setting
there is a unique structure-preserving map
$\mathcal{O}(\gamma) \to X$ sending $1 \mapsto x_0$ and intertwining
$\cdot \gamma$ with $f$.
\end{theorem}

The construction is standard once initiality is invoked. The point is
that the comparison operator does not provide $\mathbb{N}$
\emph{additionally}: it provides exactly one Peano-shaped structure up
to unique isomorphism, namely the orbit. The orbit and the standard
inductive natural-number object are, up to unique isomorphism, the
\emph{same} object.

\begin{remark}[Generator independence]
\label{rem:gen-indep}
For any two non-trivial generators $\gamma, \gamma'$, the orbits
$\mathcal{O}(\gamma)$ and $\mathcal{O}(\gamma')$ are canonically
isomorphic as Peano structures via the unique map sending
$\gamma \mapsto \gamma'$. The choice of generator does not change the
arithmetic content; it only relabels the step.
\end{remark}

\paragraph{What this section earns.} It removes the strongest form of
the import objection. The recovered natural-number object is not an
arbitrary copy of $\mathbb{N}$. It is the unique-up-to-unique-isomorphism
initial Peano structure in the comparison-operator setting. The
\emph{naming} of this object as a Lean inductive type is implementation
overhead; the \emph{mathematical} object is determined by the
comparison-operator structure alone.

\paragraph{What this section does not earn.} It does not yet show that
the same arithmetic content arises in another realization of the Law
of Logic. That is the content of the universal forcing program
(\S\ref{sec:program}).

\section{Bridge to the Universal Forcing Theorem}
\label{sec:program}

The recovery proved in this paper is one realization of a stronger
thesis: the arithmetic content of the Law of Logic is invariant
across all admissible realizations, not only the continuous
positive-ratio one. The full statement, the abstract framework, and
the survey of further realizations are the subject of the companion
paper~\cite{Washburn2026UniversalForcing}. We summarise the bridge
here so that the present paper stands on its own.

\paragraph{Abstract realization.} A \emph{Law-of-Logic realization} is
a carrier $K$ together with a composition $\cdot_K$, a symmetry
action, and a comparison operator $C : K \times K \to \mathrm{Cost}$
satisfying the operator-form analogues of the six conditions used in
\S\ref{sec:background}. The continuous positive-ratio realization
($K = \mathbb{R}_{>0}$, $\cdot_K$ multiplication, $\mathrm{Cost} =
\mathbb{R}$) is the case treated in this paper.

\paragraph{Forced arithmetic of a realization.} For any realization
$R$ with non-trivial generator $\gamma$ and identity $e_K$, define
$\mathrm{Arith}(R)$ to be the smallest $K$-subset containing $e_K$ and
closed under $\cdot \gamma$, equipped with $e_K$ as base point and
$\cdot \gamma$ as successor. The construction generalises
\S\ref{sec:inductive} verbatim. By the same argument as
Theorem~\ref{thm:initiality}, $\mathrm{Arith}(R)$ is the initial
pointed-endomap algebra inside $R$, and therefore is a natural-number
object in the Lawvere sense.

\paragraph{Universal Forcing Theorem.} The companion paper proves: for
any two Law-of-Logic realizations $R$ and $S$, the forced arithmetic
objects $\mathrm{Arith}(R)$ and $\mathrm{Arith}(S)$ are canonically
isomorphic as Peano structures. The proof uses only the universal
property of $\mathrm{Arith}(R)$ inside $R$ and of $\mathrm{Arith}(S)$
inside $S$, plus the standard initiality-uniqueness argument: the
unique structure-preserving maps in each direction compose to identity
on each side.

\paragraph{Realizations covered.} The companion paper discharges the
mathematical realizations (continuous positive-ratio, discrete
propositional, modular cyclic, order-theoretic, Lawvere-style
categorical) and presents a survey of motivating extra-mathematical
realizations (music, narrative, ethics, biology, metaphysical
ground).

\paragraph{Spacetime: time as the canonical iteration object.} The
recovered $\mathbb{N}_{\mathrm{Logic}}$ of \S\ref{sec:inductive} is
upstream-essential for the spacetime story in the framework's master
synthesis paper~\cite{Washburn2026RealityFromOneDistinction}. That
paper identifies time with the canonical orbit of the recognition
operator (``time as the forced orbit''). The identification is
honest only if the orbit is canonical, which requires the forced
arithmetic of the Law of Logic to be uniquely determined. Without
the present paper, the time-as-orbit identification would be a
structural analogy: time happens to look like the iteration object
that the Law of Logic forces, and the framework asserts the
identification on those grounds. With the present paper, the
identification is rigorous: the Law of Logic forces the comparison
operator (rigidity theorem of \cite{Washburn2026LogicFE}), the
comparison operator forces the arithmetic orbit
(\S\ref{sec:inductive}), the orbit IS the canonical iteration
object up to canonical isomorphism (Universal Forcing Theorem of
\cite{Washburn2026UniversalForcing}), and time is then identified
with that iteration object via the universal-property characterisation
of the natural-number object. The downstream spacetime derivation
(spatial dimension three from linking, Lorentzian $(1, 3)$ signature,
the eight-tick ledger kernel, the constants $c, \hbar, G$ as
$\varphi$-powers) lives in
\cite{Washburn2026RealityFromOneDistinction} and uses time-as-orbit
as a load-bearing identification. The present paper supplies the
algebraic precursor that makes the identification rigorous rather
than analogical.

\paragraph{What this implies for the present recovery.} The recovered
$\mathbb{N}_{\mathrm{Logic}}$ of \S\ref{sec:inductive} is not a
special feature of the continuous positive-ratio carrier; it is the
natural-number object that any admissible realization of the Law of
Logic forces, and the recovered Peano structure here is canonically
isomorphic to the Peano structure forced in any other realization.
The number tower constructed in
\S\ref{sec:integers}--\S\ref{sec:tower} is built on top of
$\mathbb{N}_{\mathrm{Logic}}$ by standard algebraic and completion
machinery; the upper layers therefore inherit the
realization-invariance of their Peano core. The analytic substrate
transports through Mathlib remain a separate honest reuse and do not
become realization-invariant by this argument.

\section{Current Axiom Audit}
\label{sec:audit}

The formal audit has improved since the first draft of this paper.
The Peano core of $\mathbb{N}_{\mathrm{Logic}}$ remains free of
domain-specific axioms. Quotient and completion layers use the standard
Lean kernel assumptions needed by the ambient proof assistant
(\texttt{propext}, \texttt{Classical.choice}, and \texttt{Quot.sound})
where quotients, rationals, completions, and analytic transports are
used.

The continuous-combiner extension of the companion functional-equation
paper has also been sharpened. The H-side Aczél--Kannappan classification
is now theorem-backed. The former opaque continuous-combiner input has
been split into theorem-backed algebra plus four narrow residual analysis
inputs:
\[
\begin{aligned}
&\texttt{continuous\_combiner\_mollifier\_finite\_smoothness},\\
&\texttt{continuous\_combiner\_finite\_smoothness\_to\_top},\\
&\texttt{continuous\_combiner\_second\_derivative\_identity},\\
&\texttt{continuous\_combiner\_psi\_affine\_completion}.
\end{aligned}
\]
The first two are the remaining smoothness-bootstrap content for
$G(t)=F(e^t)$; the latter two are the remaining Stetkær/Aczél-style
combiner argument. This paper does not need those residuals for the
Peano-layer recovery, but they matter for the stronger headline that
continuity alone replaces polynomial closure in the Law-of-Logic
translation theorem.

\section{Discussion}

\subsection{What is and is not claimed}

\paragraph{Claimed in this installment.} On the continuous positive-ratio
realization of the Law of Logic, the natural numbers and the integers
are canonically picked out by the comparison-operator structure, with
the universal-property characterization of \S\ref{sec:universal}. The
rationals, reals, and complex carrier are then recovered by standard
algebraic and completion constructions transported through the lower
layers. The Peano core has no domain-specific axiomatic content.

\paragraph{Claimed in the broader program.} The same arithmetic
structure is canonically forced in every admissible realization of
the Law of Logic. This is the Universal Forcing Theorem, stated and
proved in the companion paper~\cite{Washburn2026UniversalForcing}.
The present paper supplies the continuous positive-ratio realization,
which is the first mathematically dense input to that theorem.

\paragraph{Not claimed.} The recovery does not displace ZFC, ETCS, or
any other foundation of mathematics. It does not rederive analysis or
complex analysis without using the existing analytic substrate. The
proof assistant's inductive-type, quotient, and completion machinery
is reused as honest metatheoretic baseline, not as a foundational
claim.

\subsection{Comparison with set-theoretic Peano}

The standard set-theoretic Peano construction (von Neumann ordinals)
defines $0 := \emptyset$ and $n + 1 := n \cup \{n\}$, then proves
that the resulting set satisfies the Peano axioms inside ZFC. The
construction is correct; its content is that ZFC suffices to
construct $\mathbb{N}$. The construction does not exhibit $\mathbb{N}$
as a forced consequence of any classical principle outside ZFC. The
present paper exhibits $\mathbb{N}$ as a forced consequence of one
classical principle outside set theory: the Law of Logic on continuous
positive-ratio comparisons.

\subsection{Comparison with categorical Peano}

The Lawvere natural-number object is a universal property in a
topos: a triple $(N, 0, s)$ initial among triples $(X, x_0, f)$ in the
topos. The categorical presentation is correct; its content is that
$\mathbb{N}$ has a universal-property characterization. It does not
exhibit $\mathbb{N}$ as a forced consequence of any classical principle.
The present paper does, on the restricted comparison-operator setting.

\subsection{Comparison with the Curry-Howard direct construction}

In intensional type theory, $\mathbb{N}$ is an inductive type with two
constructors. What is novel here is not the inductive type but
the identification of its constructors with the orbit structure
forced by the Law of Logic: the inductive type is not arbitrary; it
is exactly what the comparison-operator structure forces.

\subsection{Foundational implications}

The recovery answers a question raised in \cite{Washburn2026LogicFE}:
whether the comparison-operator characterization extends to arithmetic.
On the present setting, it does. The Law of Logic canonically picks out
both the unique calibrated cost function and the inductive structure of
the natural numbers, with no additional axiomatic content. This makes
precise, on this setting, the philosophical claim that classical logic
and elementary arithmetic are not independent commitments but two
consequences of a single forced functional equation. The full
foundational claim, that the same arithmetic content is forced in
\emph{every} admissible realization of the Law of Logic, is the target
of the program in \S\ref{sec:program}; this paper supplies the first
realization on which the universal theorem will rest.

\section{Outlook}

\subsection{Rationals, reals, and complex analysis}

The Grothendieck completion construction now extends through the field
of fractions, completion, and complex carrier. The next mathematical
task is not to define another copy of $\mathbb{C}$, but to make the
analytic compatibility layer richer: holomorphy, contour integration,
Mellin transforms, and zeta-function identities should be stated once
over Mathlib's analytic substrate and transported back to
$\mathbb{C}_{\mathrm{Logic}}$ only where this adds mathematical content.

\subsection{Recovered analytic constants and J-cost}

The recovered real line now carries transported definitions of
$\sqrt{\cdot}$, $\exp$, $\log$, real powers, $\pi$, trigonometric and
hyperbolic functions. The RS constants then lift to recovered reals,
and the formal transport proves that each lifted constant agrees with
the previously-used real-valued constant. The canonical reciprocal
cost also lifts to recovered reals, with normalization, reciprocal
symmetry, non-negativity, square form, and zero iff identity on the
positive domain.

\subsection{Discrete propositional logic and other realizations}

The Law of Logic on continuous positive-ratio comparisons is one
realization; the discrete propositional encoding is another, and the
recovery takes a different shape there: the carrier collapses to a
finite image while the iteration object remains the same
natural-number object. The companion paper~\cite{Washburn2026UniversalForcing}
treats the discrete propositional, modular cyclic, order-theoretic,
and Lawvere-style categorical realizations. The translation theorem
of \cite{Washburn2026LogicFE} on continuous comparisons does not
generalize automatically to the discrete case; what generalises is the
universal-property characterisation of the iteration object, which is
all the Universal Forcing Theorem requires.

\subsection{Number theory and the Erd\H{o}s-Straus chain}

The recovery makes the natural, integer, rational, and real number
objects available to the existing number-theory pipeline. The
Erd\H{o}s-Straus phase-budget chain and the conditional Riemann
Hypothesis line both operate on ordinary arithmetic objects; under
the recovery, their arithmetic inputs can be re-stated in the
recovered types. The analytic zeta layer remains a separate analytic
bridge; only the arithmetic ledger hook is replaced by the recovered
prime-ledger certificate.

\subsection{Downstream axiom elimination}

Recent formal work has narrowed the downstream mathematical axiom surface.
An inconsistent legacy bounded-cost axiom in the Riemann-Hypothesis
Euler route has been replaced by an explicit deprecated proposition
which, if assumed, contradicts a proved unboundedness theorem. Classical
inputs such as Bernstein's inequality, Pick-certificate soundness, Chen,
Helfgott, Zhang--Maynard, and ABC/IUT are now recorded as sharply named
residual estimates or explicit external packages rather than silent
global axioms. This downstream cleanup is not part of the arithmetic
recovery theorem, but it clarifies which mathematical content remains
external when the recovered number tower is used in later number-theory
applications.

\appendix

\section{Formal implementation}
\label{app:formal}

The formal development is maintained at
\[
\text{\url{https://github.com/jonwashburn/recognition-science}}.
\]
The table below records the theorem surface used for the verification.
The main text deliberately treats these as implementation details; the
mathematical claim is the recovery chain itself.

\begingroup
\small
\sloppy
\begin{itemize}[leftmargin=2em]
\item \textbf{Natural numbers}: \texttt{Foundation.ArithmeticFromLogic};
      \texttt{equivNat}, \texttt{toNat\_add}, \texttt{toNat\_mul}.
\item \textbf{Integers}: \texttt{Foundation.IntegersFromLogic};
      \texttt{equivInt}, \texttt{eq\_iff\_toInt\_eq}.
\item \textbf{Rationals}: \texttt{Foundation.RationalsFromLogic};
      \texttt{equivRat}, \texttt{eq\_iff\_toRat\_eq}.
\item \textbf{Reals}: \texttt{Foundation.RealsFromLogic};
      \texttt{equivReal}, \texttt{bourbaki\_complete}.
\item \textbf{Transcendentals}: \texttt{Foundation.LogicRealTranscendentals};
      transported exponential, logarithmic, trigonometric, and hyperbolic APIs.
\item \textbf{Constants}: \texttt{Foundation.LogicRealConstants};
      transported RS constants and bounds.
\item \textbf{J-cost}: \texttt{Cost.JcostLogic};
      recovered-real normalization, square form, non-negativity, and zero
      iff identity.
\item \textbf{Law of Logic}: \texttt{Foundation.LogicAsFunctionalEquationLogic};
      recovered-real RCL transport.
\item \textbf{Complex carrier}: \texttt{Foundation.ComplexFromLogic};
      \texttt{equivComplex}, coordinate transport, and algebra transport.
\item \textbf{Tower audit}: \texttt{Foundation.RecoveredTowerAxiomAudit};
      audit surface for \texttt{LogicNat} through \texttt{LogicComplex}.
\item \textbf{Continuous combiner}: \texttt{Foundation.GeneralizedDAlembert};
      assembled continuous-combiner bilinear classification with four named
      residual analysis inputs.
\item \textbf{Universal forcing}: \texttt{Foundation.LogicRealization},
      \texttt{Foundation.ArithmeticOf}, and
      \texttt{Foundation.UniversalForcing}; realization interface,
      forced arithmetic object, canonical arithmetic invariance,
      \texttt{UniversalForcingCert}, and
      \texttt{forced\_arithmetic\_surfaces\_equivalent}.
\item \textbf{Further realizations}: \texttt{Foundation.DiscreteLogicRealization},
      \texttt{Foundation.ModularLogicRealization},
      \texttt{Foundation.OrderedLogicRealization},
      \texttt{Foundation.CategoricalLogicRealization}, and
      \texttt{Foundation.PhysicsLogicRealization}; Boolean, periodic modular,
      ordered faithful, Lawvere-style, and identity-tick realization hooks.
\item \textbf{Universal audit}: \texttt{Foundation.UniversalForcingAudit};
      audit surface for realization invariance and Peano surface.
\item \textbf{Prime ledger}: \texttt{NumberTheory.LogicPrimeLedgerAtom};
      recovered prime-ledger certificate.
\item \textbf{Erd\H{o}s-Straus}: \texttt{NumberTheory.LogicPhaseBudgetBridge};
      recovered-arithmetic phase-budget bridge.
\item \textbf{RH wrapper}: \texttt{NumberTheory.LogicRH\_From\_RCL};
      recovered-prime RH wrapper with analytic zeta bridge left explicit.
\end{itemize}
\endgroup

The current audit separates three classes of dependency:
\begin{itemize}
\item the Peano core, which is free of domain-specific axioms;
\item quotient/completion/analytic transport layers, which use the
standard Lean kernel and Mathlib infrastructure;
\item downstream number-theory applications, where external analytic
or classical inputs are named explicitly.
\end{itemize}

\section*{Acknowledgements}

We thank the team at the Recognition Physics Institute and the broader
formal mathematics community.

\begin{thebibliography}{99}

\bibitem{Washburn2026LogicFE}
Washburn, J.
\emph{A Functional-Equation Encoding of Logical Consistency on
Continuous Positive-Ratio Comparisons: A Conditional Rigidity Theorem.}
Manuscript, April 2026.

\bibitem{Washburn2026UniversalForcing}
Washburn, J.
\emph{The Universal Forcing Meta-Theorem:
One Law of Logic, One Arithmetic, Across All Realizations.}
Manuscript, April 2026.

\bibitem{Washburn2026RealityFromOneDistinction}
Washburn, J.
\emph{Reality from One Distinction:
A Forcing Chain for Spacetime, Time, and the Fundamental Constants.}
Manuscript, April 2026.

\bibitem{WashburnZlatanovicAllahyarov2026}
Washburn, J., Zlatanović, M., Allahyarov, E.
\emph{The d'Alembert Inevitability Theorem.}
Mathematics (MDPI), to appear, 2026.

\bibitem{WashburnZlatanovic2026}
Washburn, J., Zlatanović, M.
\emph{Uniqueness of the Canonical Reciprocal Cost.}
Mathematics 14 (2026), 935.

\bibitem{Aczel1966}
Aczél, J.
\emph{Lectures on Functional Equations and Their Applications.}
Academic Press, 1966.

\bibitem{Lawvere1964}
Lawvere, F. W.
\emph{An elementary theory of the category of sets.}
Proc. Natl. Acad. Sci. USA 52:1506--1511, 1964.

\end{thebibliography}

\end{document}
